\section{Local harmonic oscillations}\label{sec:local_harmonic_oscillations}
One predictive feature of FNS dynamics is the local oscillations in sampling within confined potentials. We can quantify such high-frequency oscillations under FNS dynamics by considering the deterministic part of our FNS equations of motion in \cref{eq:fns_system}. Considering a Gaussian probability distribution of fixation positions $\pi(\textbf{x}) \propto \exp \left( - \frac{\| \textbf{x} \|^2}{2 \sigma_{\text{fix}}^2} \right)$ and thus logarithmic gradients $\nabla \ln \pi(\textbf{x}) = -\frac{\textbf{x}}{\sigma_{\text{fix}}^2}$. Using this as our saliency landscape $\pi(\textbf{x})$ without considering perturbation to the explored potential from the slowly updating kernel dynamics, our equations of motion are
\begin{equation}
\begin{aligned}
    d\textbf{x}_t &= -\gamma c_{\alpha} \frac{\textbf{x}_t}{\sigma_{\text{fix}}^2} dt + \beta \textbf{p}_t dt + \gamma^{1/\alpha} dL_t^{\alpha}, \\
    d\textbf{p}_t &= -\beta c_{\alpha} \frac{\textbf{x}_t}{\sigma_{\text{fix}}^2} dt
\end{aligned}
\end{equation}
where we ignore the stochastic noise driving this equation, only to consider the coupled ODE which describes the deterministic part of $\textbf{x} = \textbf{x}_t$ and $ \textbf{p} = \textbf{p}_t$, where the loss of subscript is typical when moving from an SDE to an ODE
\begin{equation}
\begin{aligned}
    \frac{d \textbf{x}}{dt} &= -\gamma c_{\alpha} \frac{\textbf{x}}{\sigma_{\text{fix}}^2} + \beta \textbf{p}, \\
    \frac{d\textbf{p}}{dt} &= -\beta c_{\alpha} \frac{\textbf{x}}{\sigma_{\text{fix}}^2} 
\end{aligned}
\end{equation}
Assuming an underdamped regime where the system is periodic, corresponding to parameter range $\beta > \gamma \sqrt{c_{\alpha}} / 2 \sigma_{\text{fix}}$, the resonant ordinary frequency $f = \omega / 2 \pi$ is given by
\begin{equation}\label{eq:resonant_frequency}
    f = \frac{1}{2 \pi \sigma_{\text{fix}}} \sqrt{\beta^2 c_{\alpha} - \frac{\gamma^2 c_{\alpha}^2}{4\sigma_{\text{fix}}^2}}
\end{equation}
% The parameters $\beta$ and $\gamma$ drive the high-frequency dynamics observed during fixation intervals. $\beta = 1000$ was explicitly tuned to match the empirical microtremor fixation frequency at 65 Hz with a standard deviation of 3 pixels. Correspondingly, $\gamma = 1000$ was selected to be of a similar order of magnitude to $\beta$ to dampen these strong 65 Hz oscillations appropriately. The interaction of these parameters successfully replicates a prominent gamma-band peak centred at 65 Hz, matching the empirical eye-velocity power spectral density during stable fixations.
